% chapters/15/contoh.tex
\section{Dua Contoh Terapan}
\label{sec:contoh-bab-15}

\subsection{Contoh Pertama: Menghitung Kekuatan Bunyi}

\begin{contohbox}
Seorang mahasiswa mengeluh bahwa bunyi lompatnya ``terlalu pelan'', padahal
nilai \code{sfxVolume} sudah dinaikkan dari 0,5 menjadi 1,0. Musik latarnya
memakai \code{musicVolume} bernilai 0,5.
\end{contohbox}

Nilai pada Unity bukan persentase, melainkan pengali kekuatan secara lurus.
Karena itu, nilai 0,5 berarti kira-kira $-6$ dB, sedangkan nilai 1,0 berarti
$0$ dB. Dengan kedua tingkat \istilah{AudioMixer} masih pada nilai bawaannya,
keadaan sebelum perbaikan adalah:

\[
\text{musik} = -6 \text{ dB}, \qquad \text{efek} = -6 \text{ dB}.
\]

Keduanya sama kuat, dan di situlah masalahnya. Pada kekuatan yang sama, telinga
mendahulukan bunyi yang berjalan terus-menerus, sehingga setiap bunyi lompatan
tertutup oleh musik yang tidak pernah berhenti.

Setelah \code{sfxVolume} dinaikkan menjadi 1,0, keadaannya menjadi:

\[
\text{musik} = -6 \text{ dB}, \qquad \text{efek} = 0 \text{ dB},
\]

yaitu selisih 6 dB, atau efek suara dua kali lebih kuat daripada musik. Selisih
sebesar itu ternyata belum cukup. Di sinilah pengalaman mahasiswa itu berguna.
Kenaikan yang terdengar besar pada angkanya hanya menghasilkan perbedaan yang
kecil pada pendengaran, sebab 6 dB berarti dua kali kekuatan, bukan dua kali
kenyaringan yang terdengar.

Perbaikan yang benar adalah menurunkan musiknya, bukan menaikkan efek suaranya
lebih jauh:

\[
\text{musik} = -12 \text{ dB}, \qquad \text{efek} = 0 \text{ dB}.
\]

Selisih 12 dB, atau sekitar empat kali kekuatan, cukup membuat bunyi lompatan
terdengar jelas di atas musik yang sedang berjalan.

Kesimpulan yang berlaku umum: \textbf{menyeimbangkan bunyi lebih sering berarti
menurunkan yang terlalu kuat daripada menaikkan yang terlalu pelan}, dan
selisih sekitar 10 dB antara efek suara dan musik latar adalah titik awal yang
baik.

\subsection{Contoh Kedua: Menelusuri Bunyi yang Tidak Terdengar}

\begin{contohbox}
Seorang mahasiswa melaporkan bahwa bunyi mengambil benda tidak pernah
terdengar, sedangkan bunyi lompat dan bunyi terluka berjalan dengan baik.
Program berjalan tanpa pesan kesalahan apa pun.
\end{contohbox}

Penelusurannya dilakukan menurut urutan berikut.

\begin{langkah}
  \item \textbf{Periksa daftar peristiwa.} Apakah pengambilan benda memang
        sudah masuk daftar peristiwa yang memperoleh bunyi? Ternyata sudah,
        sehingga masalahnya bukan pada rencananya.
  \item \textbf{Periksa \code{Tag} pada bendanya.} Skrip memeriksa
        \code{CompareTag("Pickup")} sebelum memanggil bunyinya. Ternyata
        benda yang dapat diambil masih bertanda \istilah{Untagged}, sehingga
        pemeriksaan itu selalu gagal dan seluruh isi bloknya dilewati,
        termasuk penambahan skor.
  \item \textbf{Periksa berkasnya.} Setelah tanda itu diperbaiki, bunyinya
        masih belum terdengar. Berkasnya ternyata belum diseret ke peubah
        \code{pickupSound} pada \istilah{Inspector}.
  \item \textbf{Perbaiki keduanya, lalu jalankan kembali.} Bunyinya terdengar,
        dan skornya bertambah.
\end{langkah}

Tiga hal dapat dipetik dari penelusuran itu. Pertama, kesalahan yang sama dapat
menyebabkan dua gejala sekaligus; di sini skornya juga tidak bertambah, dan
gejala kedua itulah yang seharusnya lebih dulu menarik perhatian. Kedua,
pemeriksaan yang gagal secara diam-diam adalah penyebab yang paling sering
terjadi pada peristiwa yang dipicu sentuhan. Ketiga, rujukan yang belum
dihubungkan pada \istilah{Inspector} bernilai kosong, dan skrip pada
Daftar~\ref{lst:ch15-pemicu} melewatinya tanpa keluhan --- sebab itulah
pemeriksaan \code{if (manajer != null)} dipasang di mana-mana, dan sebab itulah
pula pemasangan yang belum lengkap perlu diperiksa dengan sengaja.
